\documentclass[11pt,a4paper]{article}
\usepackage[margin=24mm]{geometry}
\usepackage[T1]{fontenc}
\usepackage{lmodern}
\usepackage{amsmath,amssymb,graphicx,booktabs,tabularx,array,placeins}
\usepackage[font=small,labelfont=bf]{caption}
\usepackage{hyperref}
\hypersetup{hidelinks,pdftitle={Accelerating integral decomposition with Kira, machine learning, and exact reductions},pdfauthor={Rushil Saraswat}}
\setlength{\parindent}{1em}
\setlength{\parskip}{3pt}
\setlength{\emergencystretch}{2em}
\newcommand{\Q}{\mathbb Q}
\newcolumntype{Y}{>{\raggedright\arraybackslash}X}
\title{Accelerating integral decomposition with Kira, machine learning, and exact reductions}
\author{Rushil Saraswat}
\date{October 19, 2026}
\begin{document}
\maketitle
\input{paper_abstract}

\section{Introduction}
Feynman integrals arise from integration over internal loop momenta in higher-order perturbative corrections to scattering amplitudes and particle decay amplitudes. Many integrals in a calculation share propagators and differ only in their powers or numerators. Integration-by-parts (IBP) identities relate these integrals and express them in terms of a smaller set of master integrals~\cite{Weinzierl2022,ChetyrkinTkachov}. This decomposition supplies the coefficients needed to assemble an amplitude; the evaluation of any remaining master integrals is a subsequent calculation~\cite{Weinzierl2022}.

Kira automates the generation and solution of the IBP system, and its current methods already remove unnecessary seeds and equations~\cite{KiraOriginal,Kira3}. The remaining cost depends on the integral structure and the requested targets. Generating fewer useful equations can reduce work, but a restricted system can also fail to reach the required basis and require a larger calculation. We therefore optimize complete reduction time, including prediction, algebra, and recovery. Seed-count reduction alone does not measure this objective.

Our method first identifies structures that admit an exact reduction formula. It applies those formulas to complete requests, individual targets, and recognized intermediate sectors. For the remaining supported native calculations, we compare trained cost selection with a fixed seed region before Kira generates the identities. Exact algebra determines every output coefficient. Related machine-assisted searches motivate the seed intervention~\cite{Song2025,TubeSeeding2026}, while analytic one-loop reduction motivates the formula stage~\cite{DennerDittmaier2006,Hu2025}. We derive the particular formulas and cost objective used here in the appendices.

The present extension evaluates vacuum factors with repeated masses and raised or negative denominator powers by exact partial fractions. It also returns explicit Gamma-function values for the previously reduced connected one-mass trees. For native work, we train a new cost selector for sparse and complete seed regions, including the time lost to unsuccessful attempts. We compare the integrated method with the preceding application and with a fixed choice. These controls separate the contribution of analytic formulas from the contribution of ML. Earlier equation-selection and row-selection results remain separate, and the companion notes preserve their full experiments.

\section{The decomposition problem and Kira}
\label{sec:physics}
We first define the notation used in this paper. An integral family specifies inverse propagators, masses, and external scalar products. In dimensional regularization, we write~\cite{Weinzierl2022,KiraOriginal}
\begin{equation}
 I_{\boldsymbol a}(\boldsymbol x,d)=\int\prod_{h=1}^{L}
 \frac{d^d\ell_h}{i\pi^{d/2}}\prod_{j=1}^{N}D_j^{-a_j},
 \qquad D_j=Q_j^2-m_j^2+i0,\qquad d=4-2\epsilon.
 \label{eq:integral}
\end{equation}
Here $L$ is the number of loops, $\ell_h$ is loop momentum $h$, $N$ is the number of integral coordinates, and $Q_j$ is a specified linear combination of loop and external momenta. The quantity $m_j$ is the mass on line $j$, $\boldsymbol x$ contains squared masses and external invariants, $\epsilon$ is the dimensional regulator, and $i0$ specifies the Feynman pole prescription. A positive integer $a_j$ gives a denominator power, zero removes that factor, and a negative integer gives a numerator power. Auxiliary inverse propagators complete the coordinates for loop scalar products when the physical lines do not suffice~\cite{KiraOriginal}.

A sector consists of integrals with the same positive-index positions. Removing a positive position gives a sub-sector. A seed $\boldsymbol\sigma$ is an index vector at which we generate identities~\cite{KiraOriginal}. For a loop or external momentum $v$, the IBP relation is~\cite{ChetyrkinTkachov}
\begin{equation}
 0=\int\prod_h\frac{d^d\ell_h}{i\pi^{d/2}}
 \frac{\partial}{\partial\ell_k^\mu}
 \left(v^\mu\prod_jD_j^{-\sigma_j}\right)
 =\sum_u A_{\boldsymbol\sigma,o;u}(\boldsymbol x,d)
 I_{\boldsymbol a^{(u)}}.
 \label{eq:identity}
\end{equation}
The operator label $o$ specifies $(k,v)$, and $A_{\boldsymbol\sigma,o;u}$ is the coefficient of integral $u$. Appendix~\ref{app:derivation} expands the derivative into index shifts and explains why the coefficients are rational. Each seed supplies several operators and therefore several equations.

Kira orders integrals, uses symmetries and other valid relations, selects equations, and performs elimination and back-substitution~\cite{KiraOriginal,Kira3,Laporta}. For a target $T$, the result has the form
\begin{equation}
 T=\sum_{m=1}^{n_M}r_{Tm}(\boldsymbol x,d)M_m,
 \qquad r_{Tm}\in\Q(\boldsymbol x,d).
 \label{eq:reduction}
\end{equation}
Here $M_m$ is master integral $m$, $n_M$ is the number of masters, and $\Q(\boldsymbol x,d)$ is the field of rational functions with rational coefficients. Exact row operations give Eq.~\eqref{eq:reduction}, as derived in Appendix~\ref{app:derivation}. A remaining free integral in a restricted system does not by itself establish a master of the full family. We check completion against a reference basis. Kira can solve symbolically or use finite fields and reconstruct rational functions~\cite{Kira2}; this study leaves its reconstruction algorithm unchanged.

The three learned interventions occur at distinct points (Table~\ref{tab:places}). A seed policy changes which identities are generated. An equation policy selects seed--operator pairs. A row policy chooses among allowed elimination steps after coefficients exist. All three can change cost without predicting the physical coefficients themselves.
\begin{table}[!htbp]\centering\small
\caption{Learned decisions studied here. The last column states their use in the measured method.}
\label{tab:places}
\begin{tabularx}{\linewidth}{@{}lYY@{}}\toprule
Decision & Work affected & Use in this study\\\midrule
Seed region & Generation and solution of the IBP system & Cost model tested; validation retains a fixed region in the new native domain\\
Equation retention & Construction of selected IBP coefficients & Separately validated native filter\\
Elimination row & Intermediate expressions and row operations & Omitted from the retained application\\\bottomrule
\end{tabularx}
\end{table}

\section{Exact reductions before native elimination}
\label{sec:exact}
The exact stage uses the propagator structure to determine whether a formula applies. For independent loop components, the product measure and separated integrand factorize Eq.~\eqref{eq:integral}; Appendix~\ref{app:derivation} proves the resulting product of component reductions. Identical components share their calculation. Connected graphs require their own identities, so a high loop count alone does not select this shortcut.

Table~\ref{tab:exact-domains} summarizes the retained domains. For generic one-loop sectors, we resolve numerator scalar products along the external directions and average the perpendicular components. This yields scalar integrals whose raised powers are lowered by an exact IBP recurrence. Appendix~\ref{app:tensor-reduction} derives both steps. Their coefficients are rational in $d$, and keeping known linear dimension factors separate reduces repeated symbolic simplification. A singular matrix can prevent those particular recurrences from being used, which motivates the new two-point rule.
\begin{table}[!htbp]\centering\small
\caption{Exact identities used by the application. Bounds describe the implementation and do not classify all reducible integrals.}
\label{tab:exact-domains}
\begin{tabularx}{\linewidth}{@{}YY@{}}\toprule
Integral structure & Identity and accepted domain\\\midrule
Generic one-loop sectors & Tensor angular average and scalar IBP recurrence; rational kinematics, at most five coordinates, numerator and excess positive powers at most six; required matrices invertible\\
Zero-invariant two-point sectors & New endpoint formula; equal, unequal, or zero masses; rational kinematics, at most five coordinates, numerator and excess positive powers at most four\\
Scaleless vacuum families & Homogeneity gives zero; no masses or external momenta, up to 128 coordinates and thirteen loops\\
Vacuum factors & Partial fractions and analytic tadpoles; rational nonnegative masses, up to eight denominators per factor, and products through thirteen loops\\
Symbolic two-point masses & New scalar mass-derivative formula at $p^2=0$; individual powers at most six and total power at most eight\\
Connected one-scale trees & Successive massless bubble integration; massless two-point root or one massive vacuum root; implemented through thirteen loops\\
Independent products & Product of accepted component reductions; repeated components computed once\\\bottomrule
\end{tabularx}
\end{table}

\subsection{Two-point sectors at zero external invariant}
Let $D_1=\ell^2-x$ and $D_2=(\ell-p)^2-y$, where $x$ and $y$ are squared masses and $p^2=0$. Write $T(x)$ for the unit-power tadpole with squared mass $x$. Feynman parametrization gives
\begin{equation}
 B_{11}(x,y)=\frac{T(x)-T(y)}{x-y},\qquad
 B_{ab}(x,y)=\frac{\partial_x^{a-1}\partial_y^{b-1}B_{11}(x,y)}{(a-1)!(b-1)!}.
 \label{eq:light-main}
\end{equation}
Here $a$ and $b$ are positive propagator powers, and derivatives treat the two masses as independent before any specialization. Appendix~\ref{app:lightlike} derives this relation, its equal-mass limit, and the finite tensor formula. At equal masses, $B_{ab}(x,x)=T_{a+b}(x)$, where $T_n$ is the tadpole with denominator power $n$. A zero-mass tadpole vanishes in dimensional regularization~\cite{Weinzierl2022}.

For polynomial numerators, the same parameter shift leaves one radial denominator. Angular moments reduce the numerator to a finite sum, and a change from the Feynman parameter to the interpolated squared mass expresses each term through the two endpoint tadpoles. The result is a rational coefficient of $T(x)$ and $T(y)$, with no inverse external Gram matrix. This also applies when a generic one-loop recurrence reaches such a two-line sub-sector. Applying the identity at that intermediate step extends the complete reduction beyond requests that begin as bubbles.

The application tests the stated conditions before using a formula. It retains native reduction for unsupported targets, symbolic numerators outside the accepted domain, and recurrence limits. A request can contain both accepted and remaining targets; the application combines their exact outputs. These checks determine mathematical applicability. They do not ask a learned model to certify an identity.

\subsection{Direct vacuum evaluation}
A massless vacuum family has no dimensional scale and vanishes in dimensional regularization~\cite{Weinzierl2022}. We apply this check before any elimination, including to coupled loop momenta. Appendix~\ref{app:direct-vacuum} gives the radial scaling argument and the implemented conditions.

When every denominator in a component carries the same loop momentum, its integrand is a rational function of that momentum squared. Equal masses combine their powers, and partial fractions express the remaining poles through raised-power tadpoles. Appendix~\ref{app:direct-vacuum} derives the coefficients and evaluates each tadpole through Gamma functions. Polynomial terms and massless vacuum factors vanish in dimensional regularization~\cite{Weinzierl2022}. The application therefore returns both a master decomposition and an explicit analytic value for this domain.

This rule also handles repeated masses and polynomial numerators. It applies independently to each separated loop factor and computes identical components once. Recognition requires the absence of a momentum coupling between factors. Connected graphs use a separate identity.

\subsection{Connected trees and higher loop counts}
Successive massless bubble integrations also reduce specific connected trees to a unit-power integral of the same tree. For a massless two-point root, the coefficient is a rational function of $d$ and the external invariant. For a massive vacuum root, it is rational in $d$ and the one squared mass. Appendix~\ref{app:trees} derives these statements from Gamma-function products and their integer shifts. Equation~\eqref{eq:direct-tree-value} gives the complete analytic value of a connected vacuum tree, and the application now returns this value with its reduction. This family includes connected high-loop integrals, while independent products provide a separate high-loop test. Neither construction represents general coupled multiscale scattering.

\section{Learning the remaining Kira cost}
\label{sec:learning}
For supported native two-loop boxes, the model chooses one of three seed procedures before elimination: a tight sparse region, a wider sparse region, or every seed in the finite bounds used by the tight procedure. Appendix~\ref{app:sparse} defines the sparse inequalities. The complete-region action removes that sparse restriction while retaining the same Kira algebra and basis checks. It can avoid the cost of an unsuccessful restricted attempt. Recovery remains available if any chosen region is insufficient.

The training set contains 24 planar and nonplanar box batches, and a separate set of twelve batches selects the model. Each batch contains two to six targets with numerator rank one through four. Twenty-one features describe target counts, numerator distributions, and physical propagator incidence. They contain neither family names nor numerical kinematic profiles. Whole batches are disjoint across development, training, validation, and testing, although individual target vectors can recur. Appendix~\ref{app:direct-inputs} specifies the four sets.

We measure each training and validation action twice and train on its median complete time. This includes unsuccessful preliminary reductions and recovery. Cost-sensitive trees and bootstrap ensembles follow the recursive constructions of Refs.~\cite{CART1984,Breiman2001Forest}; Appendix~\ref{app:shared} derives the cost objective and lists the candidate settings. Validation selects a fixed control and enables a learned choice only when it improves the measured cost by at least one percent with no unverified selected results. An action that fails its exact checks receives a training penalty of at least 60 seconds.

A separate classifier estimates action failure from the training outcomes. Inputs outside the observed feature range or actions above the selected risk threshold use the fixed control, which must complete every training and validation request. A predicted score is an empirical decision aid, so the application still checks completion. This design trains against measured full cost while leaving every integral coefficient to exact algebra. A complete seed region can avoid an insufficient restricted system, but it can also increase the algebraic work. The timing objective accounts for both effects.

\subsection{Equation and row policies}
The equation-selection model ranks candidate seed--operator pairs before constructing their full coefficients. Its independently tested filter reduces complete time on its recorded workload. The current combined comparison keeps that filter off, so its benefit is reported separately. Multiplying the seed and equation speed ratios would assume an interaction that these measurements do not test.

The original next-row model scores allowed elimination steps in Kira's changing system. Its inference and selection costs exceeded the algebraic savings in the retained comparison. A later compiled cost policy reduced this overhead, but its interval still included equal time. We therefore omit learned row selection from the final application and retain Kira's exact elimination rules. Section~\ref{sec:results} reports the corresponding measurements.

\section{Comparison design}
\label{sec:design}
We freeze the formulas, learned policy, fixed action, and reserved requests before testing. The matched comparison contains 69 requests: thirteen two-mass vacuum products, 34 connected one-mass vacuum trees, three zero-invariant tensor bubbles, three ordinary one-loop tensors, and sixteen native box batches. Both vacuum grids cover every loop order from one through thirteen. The connected grid includes chains, stars, and binary trees, with coincident small-loop routings omitted. A separate set of 26 requests tests repeated-mass factors and coupled scaleless vacuum families, with one request per domain at each of the same loop orders.

An independent symbolic partial-fraction calculation supplies the vacuum-product references. Exact rational identities, convergent radial integrals, and a native tadpole recurrence check provide additional tests. Native box, zero-invariant, and ordinary tensor requests use independently prepared Kira references. The connected-tree checks compare with the preceding exact implementation and preserve the independently derived identity. Appendix~\ref{app:direct-inputs} states these distinctions.

The compared methods are the integrated application, the same application with its selected fixed native action, and the preceding published application. Every timed run starts a fresh process and includes imports, model loading, prediction, exact formulas, native equation generation and solution where used, recovery, output, and exit. We use one computational thread, three repetitions, and rotating method order on the same worker. Reference preparation and the final external coefficient comparison lie outside the timer. Native applications receive protected master indices where required; reference coefficients are used only for the external check. These times therefore assume the stated basis.

Each application has a 30-second limit. A request counts as complete only when every repetition agrees with its reference. The time ratio uses the sum of per-request median times on the matched completed subset. All incomplete requests remain in the completion counts. Appendix~\ref{app:statistics} defines this estimator and the request-resampling interval. The separate repeated-mass grid reports completion and time without claiming a speed factor against an unsupported control.

\input{direct_results}

\section{Conclusion}
\input{paper_conclusion}
\appendix
\section{IBP identities, row operations, and exactness}
\label{app:derivation}
The starting total-derivative identity follows from integration by parts in a region where the integral converges, followed by dimensional analytic continuation~\cite{ChetyrkinTkachov,Weinzierl2022}. Write $P_{\boldsymbol\sigma}=\prod_iD_i^{-\sigma_i}$ and $\partial_k=\partial/\partial\ell_k$. The product rule gives
\begin{align}
 \partial_k\!\cdot(vP_{\boldsymbol\sigma})
 &=(\partial_k\!\cdot v)P_{\boldsymbol\sigma}
 -\sum_i\sigma_i(v\!\cdot\partial_kD_i)
 D_i^{-\sigma_i-1}\prod_{j\ne i}D_j^{-\sigma_j}.
 \label{eq:product-rule}
\end{align}
For $v=\ell_h$, the first factor is $d\delta_{kh}$; for an external momentum it is zero. Completeness of the inverse-propagator coordinates permits us to write $v\!\cdot\partial_kD_i=u_i+\sum_jv_{ij}D_j$, where $u_i$ and $v_{ij}$ are coefficients fixed by the propagators and kinematics, with the invariant and mass terms included in $u_i$. Integrating Eq.~\eqref{eq:product-rule} therefore gives
\begin{equation}
 0=(\partial_k\!\cdot v)I_{\boldsymbol\sigma}
 -\sum_i\sigma_i u_i I_{\boldsymbol\sigma+\boldsymbol e_i}
 -\sum_{i,j}\sigma_i v_{ij}
 I_{\boldsymbol\sigma+\boldsymbol e_i-\boldsymbol e_j}.
 \label{eq:index-shifts-new}
\end{equation}
The unit vector $\boldsymbol e_i$ raises index $i$ by one. This derivation explains how a seed and an operator generate a row of the reduction system.

Let $q$ be a source row and $k$ a pivot column with $A_{qk}\ne0$. Solving the identity $\sum_jA_{qj}I_j=0$ for $I_k$ gives $I_k=-\sum_{j\ne k}(A_{qj}/A_{qk})I_j$. Subtracting $A_{rk}/A_{qk}$ times that zero identity from row $r$ yields
\begin{equation}
 A'_{rj}=A_{rj}-\frac{A_{rk}}{A_{qk}}A_{qj},\qquad A'_{rk}=0.
 \label{eq:row-new}
\end{equation}
Every step preserves the represented relation. Repetition until only prescribed master columns remain gives Eq.~\eqref{eq:reduction}. Removing seeds or equations preserves the validity of the retained identities, but can prevent this elimination from finishing. We therefore require reduction to the prescribed basis and exact agreement with the reference coefficients. Generic kinematics avoid singular pivot denominators; special kinematic limits require their own treatment.


For independent components, write the loop momenta as disjoint groups and the integrand as $\prod_b f_b$. In a convergence domain, repeated integration gives $I=\prod_b\int f_b$. Dimensional analytic continuation preserves that identity~\cite{Weinzierl2022}. Substituting each exact component reduction therefore gives
\begin{equation}
 I=\prod_b\left(\sum_m r_{bm}M_{bm}\right).
 \label{eq:product-reduction}
\end{equation}
Our product output retains this factored form. It avoids expanding the Cartesian product of master terms and uses the same representation in all timed methods.

\section{The seed regions and equation score}
\label{app:sparse}
For seed $\boldsymbol\sigma$, define excess positive power $D=\sum_i\max(\sigma_i-1,0)$, numerator rank $R=\sum_i\max(-\sigma_i,0)$, and positive sector $P=\{i:\sigma_i>0\}$. Let $h$ be the number of absent physical lines relative to the largest top sector, and let $D_T,R_T$ be the largest target values. The distance from the target numerator rectangles is defined by
\begin{equation}
 q(\boldsymbol\sigma,\mathcal T)=\min_{\boldsymbol a\in\mathcal T}
 \sum_{i\notin P}\max\{0,-\sigma_i-\max(0,-a_i)\}.
 \label{eq:seed-distance}
\end{equation}
Projection onto the interval from zero to each target numerator power leaves precisely this positive excess. Summing the distances and minimizing over targets gives Eq.~\eqref{eq:seed-distance}. We define the region by
\begin{align}
 q&\leq\max\{0,\lfloor\theta_0+\theta_1h+\theta_2D+\theta_6R\rfloor\},\nonumber\\
 R&\leq\max\{0,\lfloor R_T+\theta_3-\theta_4h\rfloor\},\qquad
 D\leq\max\{0,\lfloor D_T+\theta_5\rfloor\}.
 \label{eq:seed-regions}
\end{align}
These inequalities specify a search family rather than a sufficiency theorem. Native outer bounds, retained sector corners, and admissible target and master seeds also apply. The previously studied sparse choices use
\begin{align}
 \boldsymbol\theta_{\rm wide}&=(2.5,0,0,1.25,0,1,-0.25),\nonumber\\
 \boldsymbol\theta_{\rm tight}&=(1,0,0,0,1,1,0),\nonumber\\
 \boldsymbol\theta_{\rm medium}&=(1.5,0,0,0.5,0,1,0).
 \label{eq:seed-values}
\end{align}
The current model uses the wide and tight choices and a third action that retains all seeds within the tight procedure's native outer bounds. The medium region remains part of the preceding application. The native procedure retains its subsector truncation and recovery rules. The model selects one complete procedure before generating equations.

The separately tested equation classifier has seventeen features: four family and request dimensions; three seed counts; two differences from target power maxima; three sector descriptors; exact target membership and minimum index distance; the operator index; and the two target maxima. Positive labels mark equations retained by a complete reference reduction. Sampling at most 2,000 records of each class $y$ in batch $c$ gives each sampled record weight $N_{cy}/(n_{cy}N_c)$, where $N_{cy}$ and $n_{cy}$ are the native and sampled class counts and $N_c$ is the full batch count. Multiplying this weight by $n_{cy}$ preserves the class fraction $N_{cy}/N_c$. A leaf with weighted totals $W_0,W_1$ returns $p=W_1/(W_0+W_1)$, and a fixed threshold retains equations. The decision-tree construction follows Ref.~\cite{CART1984}. This empirical score does not certify sufficiency or probability calibration; the native completion check remains necessary.

\section{Measured-cost learning}
\label{app:shared}
Let $C_{ia}$ be the measured complete time of training request $i$ under action $a$. Unverified actions receive $C_{ia}\geq60$ seconds. This is a training penalty, and test failures are reported as failures. A leaf containing requests $S$ assigns action cost
\begin{equation}
 \widehat C_a(S)=\frac{1}{|S|}\sum_{i\in S}C_{ia},\qquad
 a_S=\arg\min_a\sum_{i\in S}C_{ia}.
 \label{eq:cost-leaf}
\end{equation}
The second formula follows because the first denominator does not depend on the action. A proposed split $S=S_1\cup S_2$ improves the empirical chosen cost by
\begin{equation}
 \Delta(S_1,S_2)=\min_a\sum_{i\in S}C_{ia}
 -\min_a\sum_{i\in S_1}C_{ia}
 -\min_a\sum_{i\in S_2}C_{ia}.
 \label{eq:cost-split}
\end{equation}
This is the parent cost minus the sum of child costs. We require $\Delta>0.002|S|$ seconds before splitting. This cost-sensitive objective uses the recursive tree and bootstrap-ensemble constructions of Refs.~\cite{CART1984,Breiman2001Forest}.

For each target, let $u$ and $v$ be its two nonnegative auxiliary numerator powers and let $R=u+v$. The 21 inputs are batch size; minimum, maximum, and mean $R$; number of distinct $R$ values; minimum, maximum, and mean of each of $u,v$; counts with $u=0$ and $v=0$; sums of $uv$ and $(u-v)^2$; number of physical propagators containing both loops; largest number of physical propagators incident on a loop; and counts with $R=1,2,3,4$. All are available before reduction.

We compare 1, 12, or 32 trees, maximum depths 1 through 5, minimum leaf populations 2, 4, or 8, and required fractional savings $\eta=0,0.03,0.08,0.15$. An ensemble averages leaf cost vectors. A single tree uses the complete training set; each ensemble tree samples training requests with replacement. The fixed action $b$ minimizes total validation cost among actions that completed every training and validation request. If $a_*$ has the least predicted cost, we choose it only when
\begin{equation}
 \widehat C_{a_*}+0.002 < (1-\eta)\widehat C_b.
 \label{eq:cost-margin}
\end{equation}
The allowance is two milliseconds per prediction during selection. Equation~\eqref{eq:cost-margin} is the defined minimum predicted saving, obtained by adding prediction overhead to the selected cost. Validation compares candidates by unverified count, complete cost, tree count, and depth, in that order. Deployment additionally requires at least one percent less validation cost and a change from the fixed action. We also fit a failure classifier for each action from its verified/unverified training labels. Each classifier has twelve trees of depth at most three and at least two requests per child. One tree uses all requests, and eleven use bootstrap samples. A leaf returns its failed fraction $p_f$. Binary Gini impurity for $n$ requests with $k$ failures is
\begin{equation}
 n\left[1-(k/n)^2-(1-k/n)^2\right]=\frac{2k(n-k)}{n}.
 \label{eq:failure-gini}
\end{equation}
The equality follows by expansion. Splits maximize the decrease in this quantity, following the classification-tree construction of Ref.~\cite{CART1984}. Averaging the leaf fractions gives an empirical risk score for each action, without a claim of calibrated failure probability. Validation selects a threshold from $0,0.05,0.15,0.30$ jointly with the cost policy. If the proposed score exceeds that threshold, or any feature lies outside its training range, the policy selects the fully verified fixed region. Validation requires no unverified selected actions before enabling the model. This guard uses observed evidence; a zero score is not a proof that an unseen reduction will complete. The application timer includes both predictions and any subsequent recovery.

\section{Physical test inputs}
\label{app:families-new}
All profiles below define squared quantities in a common fixed unit. The zero-invariant families use shifts $(0,p,q,r)$ and masses $(x,y,0,0)$, with
\begin{align}
 &p^2=0,\quad p\cdot q=-h-1,\quad p\cdot r=-h-2,\nonumber\\
 &q^2=-2h-7,\quad q\cdot r=-h-3,\quad r^2=-3h-13.
 \label{eq:light-input}
\end{align}
The unequal case has $(x,y)=(h+3,h+11)$, the equal case $(h+3,h+3)$, the mixed case $(h+3,0)$, and the massless case $(0,0)$. Targets have the form $(1+\delta_1,1+\delta_2,-b_1,-b_2)$, where $\delta_1+\delta_2$ is the excess positive power and $b_1+b_2$ is numerator rank.

Ordinary tensor families have $P$ physical lines, one auxiliary coordinate, shifts $q_1=0$, $q_{i+1}=p_i$, and
\begin{equation}
 m_i^2=h+11i^2\ (i\leq P),\quad m_{P+1}^2=0,\quad
 p_i^2=-(10h+17i),\quad p_i\cdot p_j=-(h+i+2j-2)\ (i<j).
 \label{eq:ordinary-input}
\end{equation}
Appendix~\ref{app:direct-inputs} specifies the current profiles, targets, and separate vacuum grids.

For the native planar box, let $p_1,p_2,p_3$ be massless, $2p_1\cdot p_2=s$, $2p_1\cdot p_3=t$, and $2p_2\cdot p_3=-s-t$. The seven physical line momenta are
\begin{equation}
 (\ell_1,\ell_1-p_1,\ell_1-p_1-p_2,\ell_2,
 \ell_2-p_1-p_2,\ell_2-p_1-p_2-p_3,\ell_1-\ell_2),
 \label{eq:planar-input}
\end{equation}
and the auxiliary momenta are $(\ell_1-p_3,\ell_2-p_1)$. The nonplanar routing follows Ref.~\cite{KiraOriginal}:
\begin{align}
 (Q_1,\ldots,Q_9)=(&\ell_1,\ell_2,p_1-\ell_1,p_2-\ell_2,
 p_1+p_2-p_3-\ell_1,\nonumber\\
 &p_1-\ell_1+\ell_2,p_1-p_3-\ell_1+\ell_2,
 \ell_1-p_3,\ell_2-p_1-p_2).
 \label{eq:nonplanar-input}
\end{align}
The first seven entries are physical and all masses vanish. Its convention gives $p_i^2=(p_1+p_2-p_3)^2=0$, $s=(p_1+p_2)^2$, and $t=(p_1-p_3)^2$. Expanding the squares gives $2p_1\cdot p_3=-t$ and $2p_2\cdot p_3=s+t$. The planar profiles use $(s,t)=(-1009-12h,-1021-14h)$ and the nonplanar profiles use $(-113-18h,-127-10h)$.

Every native target is $(1,1,1,1,1,1,1,-u,-v)$ with $u,v\geq0$ and $1\leq u+v\leq4$. The current split uses the profiles and batch counts in Appendix~\ref{app:direct-inputs}. Both graph types occur in all splits, so the test measures new batches and kinematic points within these structures.

\section{Time ratios and uncertainty}
\label{app:statistics}
Let $T_{imr}$ be complete elapsed time for request $i$, method $m$, and repetition $r$. For the set $\mathcal C$ completed by both methods in all repetitions, define
\begin{equation}
 S_m=\sum_{i\in\mathcal C}\operatorname{median}_{r=1,2,3}T_{imr},\qquad
 R=\frac{S_{\rm new}}{S_{\rm control}},\qquad
 \text{fractional saving}=1-R.
 \label{eq:time-ratio}
\end{equation}
These definitions compare identical completed requests. Dividing both times by the control time gives $R$, and subtracting gives the fractional saving. Completion counts remain separate because an unfinished reduction has no verified finite speedup. To estimate variation across requests, we draw $|\mathcal C|$ paired request indices with replacement and recompute $R$. We report the central 95\% of 4,000 such ratios, using the bootstrap construction of Ref.~\cite{Efron1979}. This interval describes the sampled request collection and does not establish performance for every integral family.

\input{focused_tensor_appendix}
\input{lightlike_appendix}
\input{focused_trees_appendix}
\section{Direct evaluation of vacuum factors}
\label{app:direct-vacuum}
We first identify a case that requires no master integral. If all masses vanish, there is no external momentum, and every active inverse propagator is the square of a homogeneous linear combination of loop momenta, the integrand has no dimensional scale. After Wick rotation, let $R$ be the common radial coordinate for all loop momenta. The measure and denominators give radial power $R^{\omega-1}$ with $\omega=Ld-2\sum_j a_j$. Split the radial integral at an arbitrary $\mu>0$ and continue its two pieces separately:
\begin{equation}
 \int_0^\mu R^{\omega-1}\,dR+
 \int_\mu^\infty R^{\omega-1}\,dR
 =\frac{\mu^\omega}{\omega}-\frac{\mu^\omega}{\omega}=0.
 \label{eq:direct-scaleless}
\end{equation}
The two pieces initially converge on opposite half-planes and their meromorphic continuations cancel. Angular singularities can be treated by analytic regularization before the same scaling argument. This is the standard scaleless-integral prescription in dimensional regularization~\cite{Weinzierl2022}. It gives the combined integral as zero; separating ultraviolet and infrared poles requires retaining the two regulated contributions separately. The implementation checks the masses and momentum homogeneity, rejects a zero denominator with positive power, and accepts up to thirteen loops and 128 coordinates. It excludes cut denominators and symbolic-index inputs. Couplings between loop momenta are permitted for this zero-integral rule.

Consider a loop momentum $k$ that occurs only in denominators $k^2-x_j+i0$, where $x_j$ is a squared mass. There is no external momentum in this component. Denominator powers can be positive, zero, or negative. Equal masses first combine by addition of their powers. With $z=k^2$, let $A_x$ be the resulting integer power at each distinct squared mass $x$. The integrand is the rational function
\begin{equation}
 F(z)=\prod_x(z-x)^{-A_x}.
 \label{eq:direct-rational}
\end{equation}
For a mass with $A_x>0$, multiply by $(z-x)^{A_x}$ and expand the remaining analytic factor at $z=x$. Its first $A_x$ Taylor coefficients give the principal part at that pole. Subtracting every principal part leaves a rational function without a finite pole, which is a polynomial $P(z)$. Thus
\begin{align}
 F(z)&=P(z)+\sum_{x:A_x>0}\sum_{r=1}^{A_x}\frac{c_{xr}}{(z-x)^r},\label{eq:direct-apart}\\
 c_{xr}&=\left.\frac{1}{(A_x-r)!}\frac{d^{A_x-r}}{dz^{A_x-r}}
 \prod_{y\ne x}(z-y)^{-A_y}\right|_{z=x}.\label{eq:direct-residue}
\end{align}
This derivation also covers negative powers, since they supply polynomial factors in the product. For distinct rational masses, the coefficients are rational numbers. The implementation obtains them by multiplying truncated Taylor series, so it avoids general symbolic differentiation. A factor with exponent $-A_y$ has coefficient
\begin{equation}
 [t^n](x-y+t)^{-A_y}=\binom{-A_y}{n}(x-y)^{-A_y-n}.
 \label{eq:direct-binomial}
\end{equation}
The binomial series proves Eq.~\eqref{eq:direct-binomial}; only $0\leq n<A_x$ is needed in Eq.~\eqref{eq:direct-residue}. Equal masses must be combined before this expansion, which avoids division by a vanishing mass difference.

Polynomial terms in $k$ and massless vacuum integrals have no scale and vanish in dimensional regularization~\cite{Weinzierl2022}. For a nonzero squared mass $x$, define the normalized tadpole
\begin{equation}
 T_r(x)=\int\frac{d^dk}{i\pi^{d/2}}(k^2-x+i0)^{-r}.
 \label{eq:direct-tadpole}
\end{equation}
To evaluate it, first take $x>0$ and a convergent Euclidean domain. The Schwinger representation follows from the defining Gamma integral after the change of variable $u=\alpha(q^2+x)$:
\begin{equation}
 (q^2+x)^{-r}=\frac{1}{\Gamma(r)}\int_0^\infty
 \alpha^{r-1}e^{-\alpha(q^2+x)}\,d\alpha.
 \label{eq:direct-schwinger}
\end{equation}
The Gaussian momentum integral is $\int d^dq\,e^{-\alpha q^2}=\pi^{d/2}\alpha^{-d/2}$. Inserting this in Eq.~\eqref{eq:direct-schwinger} and setting $u=\alpha x$ gives $x^{d/2-r}\Gamma(r-d/2)/\Gamma(r)$. Wick rotation supplies $(-1)^r$ for the denominator convention in Eq.~\eqref{eq:direct-tadpole}. Analytic continuation therefore gives
\begin{equation}
 T_r(x)=(-1)^r\frac{\Gamma(r-d/2)}{\Gamma(r)}(x-i0)^{d/2-r},\qquad
 \frac{T_r(x)}{T_1(x)}=
 \frac{\prod_{j=1}^{r-1}(d/2-j)}{(r-1)!x^{r-1}}.
 \label{eq:direct-gamma}
\end{equation}
The second equality follows by applying $\Gamma(z+1)=z\Gamma(z)$ repeatedly to the first. The empty product at $r=1$ is one. Applying Eq.~\eqref{eq:direct-gamma} term by term in Eq.~\eqref{eq:direct-apart} proves both the reduction to unit tadpoles and their full analytic evaluation. The result is a meromorphic expression in $d$; ultraviolet or infrared poles are retained. The same tadpole value also evaluates any unit tadpole returned by the retained zero-invariant tensor formulas. A remaining master with several denominators still requires its own evaluation.

Independent components multiply, as proved in Appendix~\ref{app:derivation}. The direct implementation accepts one through thirteen such loop momenta, at most eight denominators per component, individual absolute powers at most eight, and total positive power after equal-mass combination at most 48 per component. Masses are nonnegative rational numbers. Denominators with nonzero target power must carry one of those individual momenta with coefficient $+1$ or $-1$. These tests exclude couplings between components and specify the actual scope of the rule.

The connected one-mass vacuum trees have a separate analytic evaluation. Appendix~\ref{app:trees} derives their dimensionless Gamma product $C^V_{\mathcal T}$. With total physical power $A$, Wick rotation of all integer-power denominators gives
\begin{equation}
 V_{\mathcal T}=(-1)^A(m^2-i0)^{Ld/2-A}C^V_{\mathcal T}.
 \label{eq:direct-tree-value}
\end{equation}
The exponent follows either from the final radial integral or from the mass dimension of the $L$ measures and $A$ denominators. For a unit-power master, $A=P=2L-1$. The application now returns its explicit Gamma arguments and scale exponent together with the rational reduction. This exposes the analytic master value already implied by the tree derivation and avoids a separate master evaluation for these families.

\section{Inputs and checks for the direct-evaluation comparison}
\label{app:direct-inputs}
The two-mass vacuum grid contains one request at each $L=1,\ldots,13$. At loop order $L$, each independent component has squared masses $(h+2,h+7)$ with $h=23011+L$. Its four target power pairs are $(1,1),(2,1),(1,3),(2,-1)$, and each full target repeats the same pair across all $L$ components. These are products of one-loop factors. The connected vacuum grid has 34 requests: chains at $L=1,\ldots,13$, stars at $L=3,\ldots,13$, and binary trees at $L=4,\ldots,13$. Their parent maps are $\pi(j)=j-1$, $\pi(j)=0$, and $\pi(j)=\lfloor(j-1)/2\rfloor$, respectively. These starting orders omit coincident small-loop routings. The root squared mass remains symbolic. Each request has the unit-power master, root power two, final edge power two, and root power three; the one-loop request instead uses powers one through four. The profiles are $23111+L$, $23211+L$, and $23311+L$ for the three shapes and fix any auxiliary construction choices. These requests test preservation of the connected identity and the earlier symbolic single-loop dispatch.

The repeated-mass coverage grid uses four denominators per component with squared masses $(h+2,h+7,h+2,0)$, where $h=23511+L$. The four power assignments are $(1,1,2,0)$, $(2,1,1,1)$, $(1,2,0,-2)$, and $(0,0,0,2)$. This tests repeated masses, a massless pole, a polynomial numerator, and a scaleless target. It is checked through all thirteen loop orders, with no speed claim against an unfinished or unsupported native calculation.

The scaleless coverage grid has one family for each $L=1,\ldots,13$, with inverse propagators $\ell_i^2$ and $(\ell_i-\ell_j)^2$ for all $i<j$. It therefore contains $L(L+1)/2$ coordinates and coupled loop momenta for $L>1$. The three targets set all powers to one, raise the first power to two, or change the first power to minus one. Equation~\eqref{eq:direct-scaleless} supplies the reference zero. Native one- and two-loop checks test the same rule separately.

The three retained zero-invariant checks use the unequal, equal, and mixed mass choices of Eq.~\eqref{eq:light-input}, with profiles $23211,23212,23213$, numerator rank two, and one excess positive power. The three ordinary tensor checks use Eq.~\eqref{eq:ordinary-input} with $h=23311+P$ for $P=2,3,4$. Their target ranks and excess powers remain two and one.

The native inputs use Eqs.~\eqref{eq:planar-input} and \eqref{eq:nonplanar-input}. Development contains two batches per graph at profile 22307; training contains twelve per graph at 22343; validation contains six per graph at 22349; and testing contains eight per graph at 22367. Each batch has two to six distinct targets of the form stated in Appendix~\ref{app:families-new}. Whole batches are disjoint across these sets and exclude the preceding native collections. Individual target vectors can recur, and both graph types occur in training and testing.

An independent symbolic partial-fraction calculation checks the vacuum grid. Random rational-function checks verify that subtracting the proposed poles leaves only a polynomial. Convergent Euclidean radial quadrature checks the tadpole normalization, and a native IBP comparison checks its raised-power recurrence. Native box and ordinary tensor coefficients are checked against independently prepared Kira reductions. The connected tree comparison checks preservation against the preceding exact rule; its mathematical justification is the separate derivation in Appendix~\ref{app:trees}.

\bibliographystyle{unsrt}
\bibliography{paper_references}
\end{document}
